\documentclass[a4paper,12pt]{article}

\usepackage[T2A]{fontenc}
\usepackage[utf8]{inputenc}
\usepackage[russian]{babel}
\usepackage{amsmath,amssymb,mathtools}
\usepackage{PTSans}
\renewcommand{\familydefault}{\sfdefault}
\usepackage{icomma}
\usepackage[left=20mm,right=20mm,top=20mm,bottom=22mm]{geometry}
\usepackage{microtype}
\microtypesetup{nopatch=footnote}
\usepackage{xcolor}
\usepackage{tikz}
\usetikzlibrary{calc,arrows.meta,positioning,shapes.geometric}
\usepackage{tabularx,booktabs}
\usepackage{enumitem}
\usepackage{hyperref}
\usepackage{parskip}
\usepackage{fancyhdr}
\usepackage{setspace}
\usepackage{titlesec}

\definecolor{accent}{HTML}{008080}
\definecolor{darkaccent}{HTML}{004D4D}
\definecolor{textgray}{HTML}{333333}
\definecolor{softgray}{HTML}{F2F4F4}

\hypersetup{
  colorlinks=true,
  linkcolor=darkaccent,
  urlcolor=darkaccent,
  pdfauthor={Лёвин Артём Александрович},
  pdftitle={Чек-лист: десятичные и обыкновенные дроби}
}

\onehalfspacing
\setlength{\parindent}{0pt}
\setlength{\parskip}{0.55em}
\setlength{\emergencystretch}{2em}
\raggedbottom
\renewcommand{\arraystretch}{1.25}

\pagestyle{fancy}
\fancyhf{}
\fancyhead[L]{\small Ксения Васильченко}
\fancyhead[R]{\small 6 класс}
\fancyfoot[C]{\small \thepage}
\renewcommand{\headrulewidth}{0.3pt}
\setlength{\headheight}{14.5pt}
\renewcommand{\footrulewidth}{0pt}

\titleformat{\section}
  {\normalfont\Large\bfseries\color{darkaccent}}
  {}{0pt}{}

\titleformat{\subsection}
  {\normalfont\large\bfseries\color{darkaccent}}
  {}{0pt}{}

\setlist[itemize]{
  leftmargin=1.6em,
  itemsep=0.25em,
  topsep=0.25em
}

\setlist[enumerate]{
  leftmargin=1.9em,
  itemsep=0.45em,
  topsep=0.3em
}

\newcommand{\accentline}{%
  \par\noindent
  \textcolor{accent}{\rule{\linewidth}{0.5pt}}
  \par
}

\tikzset{
  flowstart/.style={
    ellipse,
    draw=darkaccent,
    line width=0.9pt,
    minimum width=3.4cm,
    minimum height=1.0cm,
    align=center,
    text width=3.0cm,
    inner sep=5pt
  },
  flowproc/.style={
    rounded corners=3pt,
    rectangle,
    draw=accent,
    line width=0.9pt,
    minimum height=1.05cm,
    align=center,
    text width=5.2cm,
    inner sep=6pt
  },
  flowcond/.style={
    diamond,
    aspect=2.0,
    draw=darkaccent,
    line width=0.9pt,
    align=center,
    text width=3.2cm,
    inner sep=2pt
  },
  flowarrow/.style={
    -{Latex[length=3mm,width=2mm]},
    line width=0.9pt,
    draw=textgray
  },
  flowlabel/.style={
    fill=white,
    inner sep=1.5pt,
    font=\small
  }
}

\begin{document}

\begin{titlepage}
\thispagestyle{empty}
\centering
\vspace*{2.5cm}

{\Large\textcolor{darkaccent}{Математика}}\\[1.0cm]

{\huge\bfseries Чек-лист}\\[0.45cm]

{\LARGE\bfseries
Десятичные и обыкновенные дроби}\\[0.55cm]

{\large
Сложение, вычитание, умножение, деление, сокращение}\\[1.4cm]

{\large Ксения Васильченко, 6 класс}\\[1.0cm]

{\large \today}\\[1.8cm]

{\large Лёвин Артём Александрович}\\[0.25cm]

{\large эксклюзивно для Ксении Васильченко}

\vfill

\begin{tikzpicture}[x=1cm,y=1cm]
  \draw[
    accent,
    line width=1.1pt
  ]
    (0,0)
    .. controls (3,0.65) and (7,-0.65) ..
    (10,0)
    .. controls (12,0.4) and (14,-0.25) ..
    (16,0.1);
\end{tikzpicture}

\vspace*{0.8cm}
\end{titlepage}

\section*{Основной чек-лист}
\accentline

\subsection*{1. Десятичные дроби}

\textbf{Сложение и вычитание.}
Числа записывают так, чтобы запятые находились строго
друг под другом. При необходимости справа дописывают нули:
например,
\[
1{,}2=1{,}20.
\]

Примеры:
\[
1{,}20+3{,}56=4{,}76,
\qquad
1{,}570-0{,}123=1{,}447.
\]

\textbf{Умножение.}
Множители удобно выравнивать по правому краю.
Сначала перемножают числа как целые, затем отделяют справа
столько цифр, сколько знаков после запятой было
в обоих множителях вместе.

\[
3{,}56\cdot1{,}2=4{,}272.
\]

\textbf{Деление на натуральное число.}
Запятая в частном появляется при переходе через запятую
в делимом:

\[
1{,}2:3=0{,}4.
\]

\textbf{Нули справа после запятой значения числа не меняют:}

\[
4{,}80=4{,}8,
\qquad
4{,}08<4{,}8.
\]

\textbf{Проверьте себя:}
\begin{itemize}
  \item при сложении и вычитании числа выровнены по запятой;
  \item при умножении верно подсчитано общее количество
        десятичных знаков;
  \item при сравнении сначала сопоставлены целые части,
        затем десятые, сотые и следующие разряды.
\end{itemize}

\subsection*{2. Обыкновенная дробь}

Обыкновенная дробь имеет вид

\[
\frac{a}{b},
\qquad
b\ne0.
\]

Число \(a\) называется \textbf{числителем},
число \(b\) --- \textbf{знаменателем}.

Для положительных дробей:

\begin{itemize}
  \item \textbf{правильная дробь}: \(a<b\),
        например \(\frac35\);
  \item \textbf{неправильная дробь}: \(a\ge b\),
        например \(\frac73\);
  \item \textbf{смешанное число}:
        \(2\frac13\) --- целая часть и правильная дробная часть.
\end{itemize}

\subsection*{3. Сложение и вычитание обыкновенных дробей}

Для дробей с разными знаменателями справедлива формула

\[
\frac{a}{b}\pm\frac{c}{d}
=
\frac{ad\pm bc}{bd},
\qquad
b\ne0,
\quad
d\ne0.
\]

На занятии этот приём использовался как правило
\textbf{«улыбочки»}:
два перекрёстных произведения образуют числитель,
произведение знаменателей образует знаменатель.

\textbf{Примеры:}

\[
\frac12+\frac35
=
\frac{1\cdot5+3\cdot2}{2\cdot5}
=
\frac{11}{10},
\]

\[
\frac35-\frac12
=
\frac{3\cdot2-1\cdot5}{5\cdot2}
=
\frac1{10},
\]

\[
\frac13-\frac17
=
\frac{7-3}{21}
=
\frac4{21}.
\]

\textbf{Ограничение метода.}
Произведение знаменателей всегда даёт общий знаменатель,
однако он может оказаться избыточно большим.
Если наименьший общий знаменатель виден сразу,
удобно использовать его.

\subsection*{4. Умножение и деление дробей}

\textbf{Умножение:}

\[
\frac{a}{b}\cdot\frac{c}{d}
=
\frac{ac}{bd},
\qquad
b\ne0,
\quad
d\ne0.
\]

Пример:

\[
\frac23\cdot\frac57
=
\frac{10}{21}.
\]

\textbf{Деление.}
Вторую дробь заменяют обратной и переходят к умножению:

\[
\frac{a}{b}:\frac{c}{d}
=
\frac{a}{b}\cdot\frac{d}{c},
\qquad
b\ne0,
\quad
d\ne0,
\quad
c\ne0.
\]

Пример:

\[
\frac45:\frac12
=
\frac45\cdot\frac21
=
\frac85.
\]

\subsection*{5. «Фантомная единица»}

Любое целое число можно записать в виде дроби
со знаменателем \(1\):

\[
5=\frac51,
\qquad
10=\frac{10}{1}.
\]

Так целое число приводится к форме,
удобной для обычных действий с дробями.

\textbf{Примеры:}

\[
\frac12+5
=
\frac12+\frac51
=
\frac{11}{2}
=
5\frac12,
\]

\[
\frac12:5
=
\frac12:\frac51
=
\frac12\cdot\frac15
=
\frac1{10}.
\]

\textbf{Короткий путь для сложения.}
При сложении положительного целого числа
и правильной дроби можно сразу получить смешанное число:

\[
10+\frac12=10\frac12.
\]

Для умножения и деления этот приём неприменим.

\subsection*{6. Смешанное число и неправильная дробь}

Чтобы перевести смешанное число в неправильную дробь,

\[
m\frac{a}{b}
=
\frac{mb+a}{b},
\qquad
b\ne0.
\]

\textbf{Примеры:}

\[
5\frac23
=
\frac{5\cdot3+2}{3}
=
\frac{17}{3},
\]

\[
2\frac13
=
\frac73.
\]

После преобразования выполняют действие
по обычным правилам:

\[
2\frac13+\frac12
=
\frac73+\frac12
=
\frac{14+3}{6}
=
\frac{17}{6}
=
2\frac56.
\]

Ещё один пример:

\[
2\frac13:5\frac27
=
\frac73:\frac{37}{7}
=
\frac73\cdot\frac7{37}
=
\frac{49}{111}.
\]

\clearpage
\thispagestyle{empty}

\section*{Алгоритм: выбор правила для действия с дробями}

\vspace{0.6cm}

\begin{center}
\begin{tikzpicture}[node distance=8mm]

  \node[flowstart] (start)
    {Начало};

  \node[
    flowproc,
    below=of start,
    text width=7.0cm
  ] (prep)
    {Подготовьте запись:
     целое \(n\to\frac n1\);
     смешанное число
     \(m\frac ab\to\frac{mb+a}{b}\)};

  \node[
    flowcond,
    below=of prep
  ] (action)
    {Какое действие?};

  \node[
    flowproc,
    text width=4.2cm
  ] (addproc)
    at ($(action)+(-5.1,-3.1)$)
    {Приведите к общему знаменателю
     и выполните \(+\) или \(-\)};

  \node[
    flowproc,
    below=of addproc,
    text width=4.2cm
  ] (addred)
    {Сократите дробь,
     если возможно};

  \node[
    flowstart,
    below=of addred
  ] (addend)
    {Готово};

  \node[
    flowproc,
    text width=4.2cm
  ] (mulproc)
    at ($(action)+(0,-3.1)$)
    {Перемножьте числители
     и знаменатели};

  \node[
    flowproc,
    below=of mulproc,
    text width=4.2cm
  ] (mulred)
    {Сократите дробь,
     если возможно};

  \node[
    flowstart,
    below=of mulred
  ] (mulend)
    {Готово};

  \node[
    flowproc,
    text width=4.2cm
  ] (flip)
    at ($(action)+(5.1,-3.1)$)
    {Замените деление умножением
     и переверните вторую дробь};

  \node[
    flowproc,
    below=of flip,
    text width=4.2cm
  ] (divmul)
    {Перемножьте
     полученные дроби};

  \node[
    flowproc,
    below=of divmul,
    text width=4.2cm
  ] (divred)
    {Сократите дробь,
     если возможно};

  \node[
    flowstart,
    below=of divred
  ] (divend)
    {Готово};

  \draw[flowarrow]
    (start) -- (prep);

  \draw[flowarrow]
    (prep) -- (action);

  \draw[flowarrow]
    (action.west)
    -|
    node[flowlabel,pos=0.25]
      {\(+\) или \(-\)}
    (addproc.north);

  \draw[flowarrow]
    (action.south)
    --
    node[flowlabel]
      {\(\times\)}
    (mulproc.north);

  \draw[flowarrow]
    (action.east)
    -|
    node[flowlabel,pos=0.25]
      {\(:\)}
    (flip.north);

  \draw[flowarrow]
    (addproc) -- (addred);

  \draw[flowarrow]
    (addred) -- (addend);

  \draw[flowarrow]
    (mulproc) -- (mulred);

  \draw[flowarrow]
    (mulred) -- (mulend);

  \draw[flowarrow]
    (flip) -- (divmul);

  \draw[flowarrow]
    (divmul) -- (divred);

  \draw[flowarrow]
    (divred) -- (divend);

\end{tikzpicture}
\end{center}

\clearpage

\section*{7. Сокращение дробей}

Дробь можно сокращать,
если числитель и знаменатель имеют общий множитель:

\[
\frac{ka}{kb}
=
\frac ab,
\qquad
k\ne0,
\quad
b\ne0.
\]

\textbf{Примеры:}

\[
\frac48
=
\frac{4\cdot1}{4\cdot2}
=
\frac12,
\qquad
\frac68
=
\frac{2\cdot3}{2\cdot4}
=
\frac34.
\]

После любого действия с дробями полезно проверить,
можно ли сократить полученный результат.

Например,

\[
\frac12+\frac34
=
\frac24+\frac34
=
\frac54.
\]

Знаменатель \(4\) здесь сразу является удобным
общим знаменателем.

\textbf{Критически важно:}
сокращают только \emph{множители}.
Отдельные слагаемые внутри суммы или разности
сокращать нельзя.

\[
\frac{2\cdot5}{2\cdot7}
=
\frac57,
\]

однако

\[
\frac{2+5}{2+7}
\]

по числу \(2\) не сокращается.

\subsection*{8. Разложение на простые множители}

При больших числителях и знаменателях
сокращение удобно начинать с разложения чисел
на простые множители.

Проверяют делимость на простые числа

\[
2,\ 3,\ 5,\ 7,\ 11,\ldots
\]

Примеры:

\[
24
=
2\cdot2\cdot2\cdot3
=
2^3\cdot3,
\]

\[
112
=
2\cdot2\cdot2\cdot2\cdot7
=
2^4\cdot7.
\]

Поэтому

\[
\frac{24}{112}
=
\frac{2^3\cdot3}{2^4\cdot7}
=
\frac{3}{2\cdot7}
=
\frac3{14}.
\]

\textbf{Признак делимости на \(2\):}
натуральное число делится на \(2\),
если его последняя цифра чётная.

Ещё один пример:

\[
54
=
2\cdot27
=
2\cdot3\cdot9
=
2\cdot3\cdot3\cdot3
=
2\cdot3^3.
\]

\clearpage
\thispagestyle{empty}

\section*{Алгоритм: разложение натурального числа на простые множители}

\vspace{0.8cm}

\begin{center}
\begin{tikzpicture}[node distance=9mm]

  \node[flowstart] (start)
    {Дано натуральное \(n>1\)};

  \node[
    flowproc,
    below=of start
  ] (setn)
    {Положите \(N:=n\)};

  \node[
    flowcond,
    below=of setn
  ] (done)
    {\(N=1\)?};

  \node[
    flowstart,
    right=36mm of done
  ] (end)
    {Все множители найдены};

  \node[
    flowproc,
    below=13mm of done,
    text width=6.0cm
  ] (findp)
    {Найдите наименьшее простое число \(p\),
     на которое \(N\) делится без остатка};

  \node[
    flowproc,
    below=of findp
  ] (writep)
    {Запишите множитель \(p\)};

  \node[
    flowproc,
    below=of writep
  ] (divide)
    {Замените \(N\) на \(N/p\)};

  \draw[flowarrow]
    (start) -- (setn);

  \draw[flowarrow]
    (setn) -- (done);

  \draw[flowarrow]
    (done)
    --
    node[flowlabel]
      {да}
    (end);

  \draw[flowarrow]
    (done)
    --
    node[flowlabel]
      {нет}
    (findp);

  \draw[flowarrow]
    (findp) -- (writep);

  \draw[flowarrow]
    (writep) -- (divide);

  \draw[flowarrow]
    (divide.west)
    -- ++(-2.0,0)
    |- (done.west);

\end{tikzpicture}
\end{center}

\clearpage

\section*{9. Итоговая памятка}

\begin{tabularx}
  {\linewidth}
  {@{}>{\bfseries}p{0.29\linewidth}X@{}}

\toprule
Ситуация & Что делать \\
\midrule

Десятичные \(+\) и \(-\)
&
Выравнивать числа по запятой;
при необходимости дописывать нули справа.
\\

Десятичное умножение
&
Перемножить числа как целые
и вернуть запятую по общему количеству
десятичных знаков.
\\

Дроби \(+\) и \(-\)
&
Привести к общему знаменателю;
правило «улыбочки» даёт общий знаменатель \(bd\).
\\

Дроби \(\times\)
&
Перемножить числители между собой
и знаменатели между собой.
\\

Дроби \(:\)
&
Перевернуть вторую дробь
и заменить деление умножением.
\\

Целое число рядом с дробью
&
Использовать
\[
n=\frac n1.
\]
\\

Смешанное число
&
Перевести в неправильную дробь:
\[
m\frac ab=\frac{mb+a}{b}.
\]
\\

Получен результат
&
Проверить возможность сокращения.
\\

Большие числа в дроби
&
Разложить на простые множители
и сократить общие множители.
\\

\bottomrule
\end{tabularx}

\subsection*{Типичные ошибки}

\begin{itemize}
  \item забыть выровнять десятичные дроби по запятой
        при сложении или вычитании;

  \item неверно поставить запятую
        после умножения десятичных дробей;

  \item складывать числители и знаменатели отдельно:
        \[
        \frac ab+\frac cd
        \ne
        \frac{a+c}{b+d};
        \]

  \item при делении дробей забыть перевернуть вторую дробь;

  \item начать действие со смешанными числами
        до перевода в неправильные дроби;

  \item сокращать слагаемые вместо множителей;

  \item забыть условие
        \[
        b\ne0
        \]
        для знаменателя
        и требование ненулевого делителя.
\end{itemize}

\clearpage

\section*{Тренировочный блок: Путь Б}

\textbf{11.09.2026}

\begin{enumerate}

  \item Вычислите:
        \[
        2{,}47+0{,}58.
        \]

  \item Вычислите:
        \[
        5{,}6-1{,}275.
        \]

  \item Вычислите и сократите:
        \[
        \frac13+\frac16.
        \]

  \item Вычислите и сократите:
        \[
        \frac34\cdot\frac25.
        \]

  \item Вычислите и сократите:
        \[
        \frac56:\frac{10}{9}.
        \]

\end{enumerate}

\textbf{12.09.2026}

\begin{enumerate}[resume]

  \item Вычислите:
        \[
        1{,}2\cdot3{,}5.
        \]

  \item Вычислите:
        \[
        2{,}4:6.
        \]

  \item Представьте ответ в виде смешанного числа:
        \[
        \frac7{10}+3.
        \]

  \item Представьте ответ в виде смешанного числа:
        \[
        5-\frac23.
        \]

  \item Вычислите:
        \[
        2\frac14+\frac16.
        \]

\end{enumerate}

\textbf{13.09.2026}

\begin{enumerate}[resume]

  \item Вычислите:
        \[
        3\frac25\cdot\frac58.
        \]

  \item Вычислите:
        \[
        4\frac12:1\frac15.
        \]

  \item Сократите дробь:
        \[
        \frac{84}{126}.
        \]

  \item Разложите число \(180\)
        на простые множители.

  \item Вычислите и сократите:
        \[
        \frac7{12}-\frac5{18}.
        \]

\end{enumerate}

\textbf{14.09.2026}

\begin{enumerate}[resume]

  \item Вычислите:
        \[
        1{,}75-0{,}986.
        \]

  \item Вычислите:
        \[
        \left(
        \frac35+\frac7{10}
        \right)
        :
        \frac{13}{20}.
        \]

  \item Вычислите:
        \[
        2\frac37-1\frac5{14}.
        \]

  \item Вычислите:
        \[
        3\frac13:
        \left(
        \frac56\cdot\frac45
        \right).
        \]

  \item Сократите дробь
        \[
        \frac{360}{504},
        \]
        предварительно разложив числитель и знаменатель
        на простые множители.

\end{enumerate}

\clearpage

\section*{Ответы к тренировочному блоку}

\begin{table}[ht]
\centering

\caption{Ответы к заданиям 1--20}

\begin{tabularx}
  {\linewidth}
  {@{}cX@{}}

\toprule
№ & Ответ \\
\midrule

1  & \(3{,}05\) \\

2  & \(4{,}325\) \\

3  & \(\frac12\) \\

4  & \(\frac3{10}\) \\

5  & \(\frac34\) \\

6  & \(4{,}2\) \\

7  & \(0{,}4\) \\

8  & \(3\frac7{10}\) \\

9  & \(4\frac13\) \\

10 & \(2\frac5{12}\) \\

11 & \(2\frac18\) \\

12 & \(3\frac34\) \\

13 & \(\frac23\) \\

14 & \(180=2^2\cdot3^2\cdot5\) \\

15 & \(\frac{11}{36}\) \\

16 & \(0{,}764\) \\

17 & \(2\) \\

18 & \(1\frac1{14}\) \\

19 & \(5\) \\

20 &
\(
360=2^3\cdot3^2\cdot5,
\quad
504=2^3\cdot3^2\cdot7,
\quad
\frac{360}{504}=\frac57
\).
\\

\bottomrule
\end{tabularx}

\end{table}

\end{document}